%%%%%%%%%%  scan idx 695  |  folio 682  %%%%%%%%%%
\SGAPrintPage{682}%
but such that
\[
 c_1(\underline{L}_1)c_1(\underline{L}_2)c_1(\underline{L}_3)\ne0
 \quad\text{in }H^6(X,\mu_2^{\otimes3}).
\]

\sgaStableTarget{sga6.e14.para.5-3}\textbf{5.3.} When $X$ is smooth and quasi-projective over a field $k$, one can likewise ask about the existence of a ring homomorphism
\[
 g:\operatorname{Gr}^\bullet(X)\longrightarrow A(X)
\]
with values in the Chow ring $A(X)$, taking Chern classes to Chern classes. One can repeat the preceding considerations, which show in particular how to construct $(-1)^{i-1}(i-1)!g^i$ directly in the absence of $g^i$ itself.

\sgaStableTarget{sga6.e14.para.5-4}\textbf{5.4.} To conclude this section, observe that essentially the same question arises when, in place of a scheme $X$, one considers a compact topological space $X$, say; one can again consider the associated graded $\operatorname{Gr}^\bullet(X)$ of $K^\bullet(X)$ endowed with its $\lambda$-filtration, and seek a ring homomorphism
\[
 \operatorname{Gr}^\bullet(X)\longrightarrow H^{2*}(X,\Z)
\]
taking Chern classes to Chern classes.

\sgaStableTarget{sga6.e14.sec.6}\section*{6. Cohomological Riemann--Roch theorem and Gysin homomorphism}

\sgaStableTarget{sga6.e14.para.6-1}\textbf{6.1.} For simplicity, we restrict ourselves here to schemes $X$ admitting an ample invertible Module; this implies that the naive $K^\bullet(X)$ and the other one can be identified. Let us begin with the Riemann--Roch formula ``with denominators'' for a proper complete-intersection morphism $f:X\to Y$. Since (for a prime $\ell$ prime to the residue characteristics of $Y$) one can define the $\ell$-adic Chern classes of a locally free Module $\underline{E}$,
\[
 c_i(\underline{E})\in H^{2i}(X,\Zell(i))
\]
(SGA 5 VII), corresponding to a homomorphism
\sgaStableTarget{sga6.e14.eq.6-1}\begin{equation}\tag{6.1}
 K^\bullet(X)\longrightarrow \Ch\bigl(H^{2*}(X,\Zell(*))\bigr),
\end{equation}
the question arises of giving a cohomological variant of the Riemann--Roch formula. The only ingredient missing to give the formula a meaning is a \uline{Gysin homomorphism}
%%%%%%%%%%  scan idx 696  |  folio 683  %%%%%%%%%%
\SGAPrintPage{683}%
\sgaStableTarget{sga6.e14.eq.6-2}\begin{equation}\tag{6.2}
 H^n(X,\underline{Q}_\ell(i))\longrightarrow H^{n-2d}(Y,\underline{Q}_\ell(i-d)),
\end{equation}
where $d$ is the virtual relative dimension of $f$. Assuming this homomorphism defined, and taking account of the Riemann--Roch formula in the form of Exp. VIII with values in $\operatorname{Gr}^\bullet(Y)_{\Q}$, one finds that the validity of the $\ell$-adic Riemann--Roch formula for the morphism $f$ is exactly equivalent to the commutativity of the square
\sgaStableTarget{sga6.e14.eq.6-3}\begin{equation}\tag{6.3}
\begin{tikzcd}[column sep=large,row sep=large]
 \operatorname{Gr}^\bullet(X)_{\Q} \arrow[r] \arrow[d] & H^{2*}(X,\underline{Q}_\ell(*)) \arrow[d] \\
 \operatorname{Gr}^\bullet(Y)_{\Q} \arrow[r] & H^{2*}(Y,\Qell(*)),
\end{tikzcd}
\end{equation}
where the horizontal arrows are defined from (6.1) by passing to associated graded groups and dividing so as to take Chern classes to Chern classes, and the vertical arrows are respectively the arrow of Exp. VIII and the Gysin homomorphism in cohomology (assumed given).

\sgaStableTarget{sga6.e14.para.6-2}\textbf{6.2.} Observe that the Gysin homomorphism is in any case defined when $Y$ is smooth over a field $k$ (SGA 5 IV), using the relative cohomological purity theorem over $k$; the same construction applies when $Y$ is excellent, regular and of characteristic zero, by the absolute cohomological purity theorem of SGA 4 XIX, and it will apply automatically to every excellent regular $Y$, without restriction on the characteristic, once the absolute cohomological purity theorem has been proved in full generality. Finally, the same construction applies, by the relative cohomological purity theorem (SGA 4 XVI 3.7), when $X$ and $Y$ are both smooth over the same scheme $S$ and $f$ is an $S$-morphism. One should conjecture that, in each of these cases in which the Gysin homomorphism is defined, diagram (6.3) is indeed commutative. The only case in which this assertion is currently verified is that in which $X$ and $Y$ are both smooth over a field $k$; it then follows from the commutativity of the diagram analogous to (6.3):
\sgaStableTarget{sga6.e14.eq.6-4}\begin{equation}\tag{6.4}
\begin{tikzcd}[column sep=large,row sep=large]
 A(X) \arrow[r] \arrow[d] & H^{2*}(X,\Zell(*)) \arrow[d] \\
 A(Y) \arrow[r] & H^{2*}(Y,\Zell(*)),
\end{tikzcd}
\end{equation}
%%%%%%%%%%  scan idx 697  |  folio 684  %%%%%%%%%%
\SGAPrintPage{684}%
involving the rings of cycle classes (mod. rational equivalence) $A(X)$ and $A(Y)$, whose commutativity is established in SGA 5 IV.

\sgaStableTarget{sga6.e14.para.6-3}\textbf{6.3.} Observe that, in each of the cases considered above in which one can define a Gysin homomorphism, one in fact obtains a construction for cohomology with coefficients in $\Zell(*)$, and not only in $\Qell(*)$. When $f$ is an \uline{immersion}, this allows the question of the validity of the Riemann--Roch formula ``without denominators'' to be posed, as explained in no.~3. When $X$ and $Y$ are both smooth over a field $k$, the commutativity of (6.4) shows that it would suffice to prove the formula without denominators in the Chow ring $A(Y)$ cf. (3.3).

\sgaStableTarget{sga6.e14.para.6-4}\textbf{6.4.} Without a regularity hypothesis on $Y$, even when $X$ and $Y$ are of finite type over the field $\C$, one no longer sees a plausible definition of a cohomological Gysin homomorphism (6.2), even for a regular immersion of codimension $1$. Can one nevertheless find one with reasonable properties, in particular making square (6.3) commute? The question is clearly connected with that of defining, for a complete intersection $X$ of codimension $i$ in $Y$, ``its cohomology class''
\[
 \gamma(X)\in H^{2i}(Y,\Zell(i)),
\]
so as to take intersections (without excess components) to cup products. This latter question has an affirmative answer at least modulo torsion, i.e. in $H^{2*}(Y,\Qell(*))$, by defining
\[
 \gamma(X)=1/(-1)^{i-1}(i-1)!\,c_i(\underline{O}_X).
\]
But I do not know whether $c_i(\underline{O}_X)$ is divisible by $(i-1)!$ in $H^{2i}(Y,\Zell(i))$. In particular, for a projective scheme $Y$ of dimension $3$ over the field of complex numbers, say, with $X$ reduced to a (singular) point of $Y$, is $c_3(\underline{O}_X)$ divisible by $2$ in $H^6(Y,\Z_2(3))$, and in particular does the Chern class vanish in $H^6(Y,\Z/2\Z)$? In the same direction, let us also mention the following question: let $Y$ be a projective (singular) scheme of dimension $4$, and $X$ a closed subscheme of dimension $1$ that is globally an intersection of three positive divisors $D_1,D_2,D_3$; consider the product of their Chern classes mod. $2$:
\[
 c_1(D_1)c_1(D_2)c_1(D_3)\in H^6(Y,\Z/2\Z);
\]
%%%%%%%%%%  scan idx 698  |  folio 685  %%%%%%%%%%
\SGAPrintPage{685}%
is this element of $H^6$ independent of the particular choice of the $D_i$ with intersection $X$?

\sgaStableTarget{sga6.e14.sec.7}\section*{7. Chern classes of perfect complexes}

\sgaStableTarget{sga6.e14.para.7-1}\textbf{7.1.} In seeking to develop a cohomological Riemann--Roch formula, with or without denominators, beyond the rather artificial hypothesis that the schemes under consideration admit an ample invertible Module, one is necessarily confronted with the question of defining, for a perfect complex of Modules $L_\bullet$ on a scheme $X$, Chern classes
\[
 c_i(L_\bullet)\in H^{2i}(X,\Zell(i)),
\]
generalizing the Chern classes of locally free Modules in SGA 5 VII. The same question also arises for general locally ringed topoi (compare [10]) and for perfect complexes on paracompact spaces ringed by continuous complex functions, working then with ordinary integral cohomology. For a compact space, the $K^\bullet(X)$ constructed with perfect complexes is known to be the same as the ``naive'' one constructed with locally free Modules (Exp. IV), and in this case one is trivially reduced to the classical definition. But even for a locally compact space the problem posed seems nontrivial, since the integral cohomology of such a space is not always the inverse limit of the cohomologies of its compact subsets.

\sgaStableTarget{sga6.e14.para.7-2}\textbf{7.2.} More precisely, for a closed subset $Y$ of $X$ such that $L_\bullet|(X-Y)$ is zero (in the sense of derived categories), i.e. such that the cohomology sheaves of $L_\bullet$ have supports in $Y$, we would like to define Chern classes \uline{with support in} $Y$
\[
 c_i^Y(L_\bullet)\in H_Y^{2i}(X,\Zell(i))\quad\text{resp.}\quad H_Y^{2i}(X,\Z),
\]
which would recover the ``ordinary'' (sic) Chern classes of 7.1 by means of the canonical homomorphism
\[
 H_Y^*(X,-)\longrightarrow H^*(X,-).
\]
%%%%%%%%%%  scan idx 699  |  folio 686  %%%%%%%%%%
\SGAPrintPage{686}%
Even when these reduce to the Chern classes of locally free Modules ($X$ a scheme with an ample invertible Module, or a compact space), these Chern classes ``with supports'' have not been constructed, and their construction will doubtless have to proceed together with that of the general Chern classes postulated in 7.1\SGAAuthorFootnote{*}{Added in October 1968. The problems raised here have been solved in the setting of \uline{Hodge cohomology} by L. ILLUSIE; cf. [13]}

It is probable that Chern classes with supports will have a role to play in proving the cohomological Riemann--Roch formula without denominators (cf. 3.5). On the other hand, they should allow an improved version of the Riemann--Roch formula (with or without denominators) by introducing families of supports into that formula.

\sgaStableTarget{sga6.e14.sec.8}\section*{8. Chow ring of regular schemes}

The problem is to define rational equivalence of cycles and to study the resulting ring of classes (in particular its functorial properties) for a regular noetherian scheme $X$ admitting an ample invertible Module, generalizing the theory known when $X$ is smooth over a field [4]. The key technical result will again be a ``moving lemma'', ensuring that a cycle can be varied within its class so as to intersect a given cycle in the correct dimension.

The problem considered is quite natural a priori, particularly since we have seen (4.5 to 4.8) that the substitutes $\operatorname{Gr}^\bullet(X)$, $\Gtop^\bullet(X)$ for the Chow ring $A(X)$ are hardly satisfactory beyond a theory ``modulo torsion''. When $X$ is obtained from a smooth projective algebraic scheme $Y$ over a field $k$ by taking the local scheme at the vertex of a projecting cone $C$ and removing the origin, one finds (admitting that $A(X)$ is indeed the direct limit of the Chow rings of the $U-(O)$, where $U$ ranges over the open neighbourhoods in $C$ of the vertex $O$) that $A(X)$ is canonically isomorphic to the quotient ring $A(Y)/\xi A(Y)$, where $\xi\in A^1(Y)$ is induced by the projective immersion of $Y$ under consideration. This shows that knowledge of the Chow ring $A(X)$ (of a nonclassical type) is equivalent to knowledge of $A(Y)/\xi A(Y)$, which might be called the ``primitive part'' of the Chow ring $A(Y)$ (in a sense suggested by the Lefschetz--Hodge cohomological theory of nonsingular projective algebraic varieties). This example shows that the ``local Chow rings'' of local rings with a singularity
